% !TeX program = xelatex
% UTF-8. Standalone: xelatex GERD.tex from this directory.
% For inclusion in a XePersian document, define \GERDInNotebook first.
% Source: Google Drive notebook/ped GERD; Nelson [22] - 369.
% Original JPEGs in filename/handwritten-page order:
% pic01.jpg = 20261016_211309871.jpg (1RC_9UDenj5OKwS0I0-pbC6uyq1yFlS3Z)
% pic02.jpg = 20261016_211323257.jpg (1Ut7hlws9xoJ4Owa_sJcCBfiCIw8-9iNl)
% pic03.jpg = 20261016_211329612.jpg (1VlfLECa04uOUjZ8o271Y7Uvdaf8SttTi)
% pic04.jpg = 20261016_211336416.jpg (1fCiogqSWYdh3cCHSV6Wyax9GtFV5b5QG)
\ifdefined\GERDInNotebook
  \let\GERDFinish\relax
\else
  \def\GERDFinish{\end{document}}
  \documentclass[12pt,a4paper]{article}
  \usepackage[margin=22mm]{geometry}
  \usepackage{amsmath,amssymb}
  \usepackage{graphicx}
  \usepackage[normalem]{ulem}
  \usepackage{xepersian}
  \IfFontExistsTF{Amiri}{\settextfont{Amiri}}{\settextfont{DejaVu Sans}}
  \setlatintextfont{DejaVu Serif}
  \setlength{\parindent}{0pt}
  \setlength{\parskip}{5pt}
  \setlength{\emergencystretch}{3em}
  \begin{document}
\fi
\providecommand{\GERDPage}[1]{\clearpage\section*{تصویر #1}}
\providecommand{\GERDUnclear}[1]{\textbf{[خوانش نامطمئن: #1]}}
\providecommand{\GERDArrow}{\lr{\(\longleftarrow\)}}
\providecommand{\GERDImagePath}{images/}
\IfFileExists{images/pic01.jpg}{}{%
  \renewcommand{\GERDImagePath}{pediatrics/ped-gastroenterology/GERD/images/}}

\begin{center}
{\Large\bfseries \lr{GERD}}\par
\lr{Nelson [22] -- 369}
\end{center}
\textbf{یادداشت رونویسی:} متن فارسی و انگلیسی، اعداد، اختصارات و حاشیه‌ها به ترتیب چهار تصویر بازنویسی شده‌اند. ارتباط فلش‌ها و شاخه‌های نمودار به صورت فهرست حفظ شده است. متن خط‌خورده جداگانه مشخص شده و خوانش نامطمئن علامت دارد. تاریخ چاپی هر چهار برگه خالی است: / /؛ شمارهٔ دست‌نویس برگه‌ها ۱ تا ۴ است. تصاویر کامل در پیوست آمده‌اند.

\section*{تصویر ۱}
\lr{GERD} شایع‌ترین اختلال مری در کودکان در تمام سنین است.

نامحتمل است که \lr{GER} فیزیولوژیک در زیر ۱ هفتگی یا بعد از ۶ ماهگی دیده شود.
\begin{itemize}
\item پیک آن ۳–۴ ماهگی است.
\end{itemize}

چه زمانی پدیده پاتولوژیک می‌شود؟ \GERDArrow{} وجود شواهد ازوفاژیت، شواهد اکستراازوفاژیال (سرفه، ویز، خشونت صدا)، نارسایی رشد.

عوامل مرتبط با \lr{LES} که دفعات رفلاکس را تعیین می‌کنند:
\begin{itemize}
\item تون ناکافی \lr{LES}، فرکانس غیرطبیعی ریلکس شدن \lr{LES}، وجود هرنی هیاتال.
\end{itemize}
طول اپیزودهای رفلاکس در شرایطی که بلع کم می‌شود افزایش می‌یابد \GERDArrow{} مثلاً حین خواب.

\subsection*{\lr{Transient LES Relaxation (TLESR)}}
\begin{itemize}
\item مکانیسم پاتوفیزیولوژی رخ دادن رفلاکس است.
\item ریلکس شدن همزمان \lr{LES} و \lr{Crura} اطراف آن.
\item مستقل از بلع رخ می‌دهد.
\item طی آن فشار \lr{LES} به حدود \lr{0--2 mmHg} بالاتر از فشار معده کاهش می‌یابد.
\item ۱۰–۶۰ ثانیه طول می‌کشد.
\item دیستنشن معده (پست پرندال، بلع هوا، تخلیهٔ غیرطبیعی معده) مهم‌ترین محرک آن است.
\end{itemize}

\subsection*{اپیدمیولوژی}
\begin{itemize}
\item رفلاکس طبیعی فیزیولوژیک در ۴ ماهگی اوج شیوعش است و در ۸۸\% موارد تا ۱۲ ماهگی رفع می‌شود و تا ۲۴ ماهگی در تقریباً ۱۰۰\% موارد رفع می‌شود.
\item \lr{Happy Spitters}: شیرخوارانی که رگورژیتیشن مکرر بدون علامت دارند.
\item علائم \lr{GERD} در کودکان بزرگ‌تر، سیر مزمن، \lr{Wax/Wan}، که در کمتر از ۵۰\% موارد به طور خودبه‌خود رفع می‌شود \GERDArrow{} بقیه تا بزرگسالی می‌مانند.
\item حتی اگر خودبه‌خود رفع شود، شواهد بافت‌شناختی ازوفاژیت باقی می‌ماند.
\end{itemize}

\GERDPage{۲}
\subsection*{تظاهرات بالینی}
\begin{itemize}
\item رفلاکس فیزیولوژیک عموماً با رگورژیتیشن تظاهر می‌یابد \GERDArrow{} خصوصاً پست پرندال.
\item شواهد ازوفاژیت \GERDArrow{} تحریک‌پذیری و عدم تمایل به شیر، خفگی.
\item \lr{FTT}.
\item علائم در تقریباً اکثر شیرخواران تا ۱۲–۲۴ ماهگی رفع می‌شود \GERDArrow{} اگر فیزیولوژیک باشد.
\end{itemize}
در سنین پره‌اسکول نیز برخی کودکان دچار رگورژیتیشن می‌شوند.

\subsection*{سندرم \lr{Sandifer}}
\lr{Refuse} کردن شیر و غذا و برگرداندن سر (\lr{Arching}).

\textbf{متن خط‌خورده در منبع:} «بروز حملات خاصی»؛ «با \lr{GERD} تفاوت است».

در زمینهٔ برخی اختلالات از جمله \lr{GERD}.

\subsection*{تظاهرات تنفسی \lr{GERD}}
تظاهرات تنفسی \lr{GERD} بر اساس سن متفاوت است:
\begin{itemize}
\item در شیرخواران به صورت آپنهٔ انسدادی یا استریدور یا \lr{Lower airway disease}.
\item رفلاکس، لارنگومالاسی یا دیسپلازی \lr{BP} را عارضه‌دار می‌کند.
\end{itemize}
اوتیت مدیا، سینوزیت، هایپرپلازی لنفوئید، خشونت صدا، ندول وکال کورد، ادم لارنکس.
\begin{itemize}
\item در کودکان بزرگ‌تر بیشتر شبیه آسم، سینوزیت، لارنژیت.
\end{itemize}

\subsection*{تشخیص}
برای اکثر بیماران با تظاهرات تیپیک، خصوصاً در سنین بالاتر، شرح حال و معاینه برای تشخیص کافی است.
\begin{itemize}
\item پرسشنامهٔ \lr{I-GERQ} شرح حال ما را استاندارد می‌کند.
\item بررسی \lr{alarm sign} ضروری است:
\par استفراغ صفراوی، استفراغ جهندهٔ مکرر، \lr{GIB}، \GERDUnclear{تاندونی}، ارگانومگالی، دیستنشن، دیسفاژی، اودینوفاژی، میکرو و ماکروسفالی، ارگانومگالی، آنورکسی، \lr{FTT}، اسهال، تب، فونتانل بالج، تشنج.
\end{itemize}
در کودک با استفراغ و دیسفاژی برای بررسی آشالازی، انترتیک وب، استنوز، هرنی هیاتال، \lr{GOO}:
\begin{itemize}
\item رادیوگرافی با کنتراست (باریم).
\item \lr{Sen, Spe} پایین در تشخیص رفلاکس دارد.
\end{itemize}

\GERDPage{۳}
\subsection*{\lr{pH monitoring}}
شرط ضروری (\lr{Sin qua non}) برای تشخیص \lr{GERD} نیست.
\begin{itemize}
\item مهم‌ترین کاربرد آن \GERDArrow{} بررسی کفایت سرکوب اسید.
\item تشخیص بیماری \lr{Atypical} (سرفه).
\end{itemize}

\subsection*{اندوسکوپی}
امکان بررسی از نظر ازوفاژیت اروزیو را به ما می‌دهد.

\subsection*{سینتی‌گرافی رادیونوکلئوتید}
بررسی آسپیراسیون در زمینهٔ رفلاکس و بررسی تخلیهٔ معده.

مانومتری مری برای تشخیص \lr{GERD} کاربردی ندارد ولی برای بررسی \lr{TLESR} مفید است.

\subsection*{برونکوسکوپی و انجام \lr{BAL}}
بررسی \lr{Silent Aspir} \GERDArrow{} بررسی کمی ماکروفاژهای \lr{lipid-laden}.
\begin{itemize}
\item تشخیص وجود پپسین در \lr{BAL}.
\end{itemize}

\subsection*{درمان امپریکال با \lr{PPI}}
به صورت کوتاه‌مدت و \lr{high dose}.
\begin{itemize}
\item اگر بیمار پاسخ نداد اپروچ تشخیص خودمان را دنبال می‌کنیم.
\end{itemize}

\subsection*{درمان \lr{GERD}}
هستهٔ اصلی درمان، محافظه‌کارانه و تغییر \lr{lifestyle} است.
\begin{itemize}
\item اصلاح الگو و روش شیردهی و غذا.
\item استفاده از غلیظ‌کننده‌های شیر \GERDArrow{} در پره‌ترم‌ها برخی از آن‌ها خطر \lr{NEC} را بالا \GERDUnclear{می‌برند؛ انتهای عبارت در حاشیهٔ تصویر قطع شده است}.
\item یک ترایال ۲ هفته‌ای با رژیم هایپوآلرژن در شیرخوار، قبل از درمان‌های دارویی.
\item عدم مواجهه با دود دخانیات.
\item \GERDUnclear{مداخلهٔ بالا} علائم را در ۲۵–۴۰\% شیرخواران کاهش می‌دهد.
\item کاهش مصرف گوجه‌فرنگی، شکلات، نعنا، آب‌میوه، نوشابه.
\item کاهش وزن در کودک با اضافه‌وزن.
\end{itemize}

پوزیشن نشسته رفلاکس شیرخوار را تشدید می‌کند. وضعیت \lr{Supine} موارد رفلاکس را نسبت به \lr{prone} تشدید می‌کند ولی چون \lr{prone} با مرگ ناگهانی همراه است، وضعیت خواب باید \lr{Supine} باشد.

پس از شیر خوردن باید \lr{flat supine} یا \lr{semi-seated pos} باشد.

\GERDPage{۴}
\subsection*{اپروچ تشخیصی بر اساس گایدلاین \lr{LA}}
سوزش سردل و/یا رگورژیتیشن بدون \lr{Alarm Sign} (علائمی که \lr{QOL} را مختل می‌کنند).
\begin{enumerate}
\item ترایال ۸ هفته‌ای با \lr{PPI}، روزانه.
\item دو شاخهٔ پاسخ به ترایال:
  \begin{itemize}
  \item رفع کامل علائم (\lr{GERD} محتمل است) \GERDArrow{} قطع \lr{PPI} \GERDArrow{} بازگشت علائم \GERDArrow{} اندوسکوپی بدون \lr{PPI} (۲–۴ هفته).
  \item رفع ناکامل علائم \GERDArrow{} اندوسکوپی بدون \lr{PPI} (۲–۴ هفته).
  \end{itemize}
\item نتیجهٔ اندوسکوپی:
  \begin{itemize}
  \item \lr{NL} یا \lr{LA} گرید \lr{A} \GERDArrow{} مانیتورینگ رفلاکس بدون درمان \GERDArrow{} مختل \GERDArrow{} تأیید تشخیص \lr{GERD}.
  \item گرید \lr{LA B/C/D} یا بارت \lr{\(>3\)} سانتی‌متر \GERDArrow{} تأیید تشخیص \lr{GERD}.
  \end{itemize}
\end{enumerate}

\subsection*{داروهای ضد اسید}
داروهای ضد اسید به سرعت اسید را خنثی کرده و علائم را رفع می‌کنند.
\begin{itemize}
\item مصرف طولانی‌مدت توصیه نمی‌شود \GERDArrow{} \lr{Al} یبوست و \lr{Mg} اسهال می‌دهد.
\end{itemize}

\subsection*{مهارکننده‌های \lr{H\textsubscript{2}} (\lr{H\textsubscript{2}RAs})}
فاموتیدین.
\begin{itemize}
\item درمان ازوفاژیت خفیف تا متوسط.
\end{itemize}
\lr{PPI} از \lr{H\textsubscript{2}RA} بهتر است در درمان ازوفاژیت شدید و اروزیو.

\subsection*{داروهای پروکینتیک}
داروهای پروکینتیک فشار \lr{LES} را \lr{\(\uparrow\)}، تخلیهٔ معده را تسریع می‌کنند ولی اثری بر \lr{TLESR} ندارند.
\begin{itemize}
\item متوکلوپراماید (آنتاگونیست \lr{D2} و \lr{5-HT\textsubscript{3}}) \GERDArrow{} نباید بیشتر از ۳ ماه استفاده شود \GERDArrow{} ریسک تاردیو دیس‌کینزی.
\item اریترومایسین.
\end{itemize}

\clearpage
\section*{پیوست: تصاویر اصلی}
\begin{center}
\includegraphics[width=\linewidth,height=0.87\textheight,keepaspectratio]{\GERDImagePath pic01.jpg}
\end{center}
\clearpage
\subsection*{\lr{pic02.jpg}}
\begin{center}
\includegraphics[width=\linewidth,height=0.87\textheight,keepaspectratio]{\GERDImagePath pic02.jpg}
\end{center}
\clearpage
\subsection*{\lr{pic03.jpg}}
\begin{center}
\includegraphics[width=\linewidth,height=0.87\textheight,keepaspectratio]{\GERDImagePath pic03.jpg}
\end{center}
\clearpage
\subsection*{\lr{pic04.jpg}}
\begin{center}
\includegraphics[width=\linewidth,height=0.87\textheight,keepaspectratio]{\GERDImagePath pic04.jpg}
\end{center}
\GERDFinish
